\documentclass[letterpaper, 12pt]{article}
\usepackage[utf8]{inputenc}
\usepackage[spanish]{babel}
\usepackage{amsmath, amssymb, amsthm}
\usepackage{geometry}
\usepackage{listings}
\usepackage{xcolor}
\usepackage{graphicx}
\usepackage{enumitem}

\geometry{left=2.5cm, right=2.5cm, top=2.5cm, bottom=2.5cm}

% Configuración para código Python
\definecolor{codegreen}{rgb}{0,0.6,0}
\definecolor{codegray}{rgb}{0.5,0.5,0.5}
\definecolor{codepurple}{rgb}{0.58,0,0.82}
\definecolor{backcolour}{rgb}{0.95,0.95,0.92}

\lstdefinestyle{mystyle}{
    backgroundcolor=\color{backcolour},   
    commentstyle=\color{codegreen},
    keywordstyle=\color{magenta},
    numberstyle=\tiny\color{codegray},
    stringstyle=\color{codepurple},
    basicstyle=\ttfamily\footnotesize,
    breakatwhitespace=false,         
    breaklines=true,                 
    captionpos=b,                    
    keepspaces=true,                 
    numbers=left,                    
    numbersep=5pt,                  
    showspaces=false,                
    showstringspaces=false,
    showtabs=false,                  
    tabsize=2
}
\lstset{style=mystyle}

\title{\textbf{Solucionario: Matemáticas Básicas}\\Parte 1}
\author{Asistente de IA}
\date{\today}

\begin{document}

\maketitle
\tableofcontents
\newpage

\section{Capítulo 1: Lógica y Conjuntos}

\subsection{1.1.11 Ejercicios (Lógica Proposicional)}

\begin{enumerate}
    \item \textbf{Determine si las siguientes f.b.f son tautologías:}
    
    \begin{enumerate}
        \item $p \to (q \to p)$ \\
        \textbf{Solución:} \textbf{Es Tautología}. Conocida como la Ley de Afirmación del Antecedente. Si $p$ es V, entonces $(q \to V)$ es V, y $V \to V$ es V. Si $p$ es F, $F \to (\dots)$ siempre es V.
        
        \item $(p \to q) \to p$ \\
        \textbf{Solución:} \textbf{No es tautología}. Si $p$ es Falso y $q$ es Verdadero: $(F \to V)$ es $V$, y luego $V \to F$ es Falso.
        
        \item $(p \lor q) \to r \iff (p \to r) \land (q \to r)$ \\
        \textbf{Solución:} \textbf{Es Tautología}. Es una equivalencia lógica estándar (prueba por casos).
        
        \item $p \to (q \to r) \iff (p \land q) \to r$ \\
        \textbf{Solución:} \textbf{Es Tautología}. Ley de Exportación.
        
        \item $p \to (p \lor q)$ \\
        \textbf{Solución:} \textbf{Es Tautología}. Ley de Adición.
        
        \item $(p \land q) \to (q \land p)$ \\
        \textbf{Solución:} \textbf{Es Tautología}. Propiedad Conmutativa.
    \end{enumerate}

    \item \textbf{Evaluación de valores de verdad:} \\
    Dadas: $p$: ``Hace frío'' (V), $q$: ``Es de noche'' (F).
    
    \begin{itemize}
        \item ``No es de noche o no hace frío'': $\neg q \lor \neg p \equiv \neg(F) \lor \neg(V) \equiv V \lor F \equiv \textbf{Verdadero}$.
        \item ``Hace frío o no es de noche'': $p \lor \neg q \equiv V \lor V \equiv \textbf{Verdadero}$.
        \item ``Ni es de noche ni hace frío'': $\neg q \land \neg p \equiv V \land F \equiv \textbf{Falso}$.
        \item ``Es de noche pero no hace frío'': $q \land \neg p \equiv F \land F \equiv \textbf{Falso}$.
    \end{itemize}

    \item \textbf{Negaciones:}
    \begin{itemize}
        \item De ``No es de noche o no hace frío'' ($\neg q \lor \neg p$):
        Negación: $\neg(\neg q \lor \neg p) \equiv q \land p$.
        En español: ``Es de noche y hace frío''.
        \item De ``Hace frío o no es de noche'' ($p \lor \neg q$):
        Negación: $\neg(p \lor \neg q) \equiv \neg p \land q$.
        En español: ``No hace frío y es de noche''.
    \end{itemize}
\end{enumerate}

\subsection{1.2.4 Ejercicios Propuestos (Conjuntos)}

\begin{enumerate}
    \item \textbf{¿Puede suceder que $A \cup B = B$?} \\
    \textbf{Solución:} Sí. Condición necesaria y suficiente: $A \subseteq B$. Si $A$ está contenido en $B$, al unirlos no se añade nada nuevo a $B$.
    
    \item \textbf{¿Puede suceder que $A \cap B = A$?} \\
    \textbf{Solución:} Sí. Condición necesaria y suficiente: $A \subseteq B$. Si todos los elementos de $A$ están en $B$, su intersección es todo $A$.
    
    \item \textbf{¿Puede suceder que $A - B = \emptyset$?} \\
    \textbf{Solución:} Sí. Ocurre si y solo si $A \subseteq B$. No hay elementos en $A$ que sean "exclusivos" de $A$.

    \item \textbf{Diagramas de Venn (Descripción):}
    \begin{itemize}
        \item a) $A \subseteq B \subseteq C$: Círculos concéntricos, $A$ dentro de $B$, $B$ dentro de $C$.
        \item b) $(A \cap C = \emptyset) \land (B \cap C = \emptyset)$: $C$ está separado (disjunto) de $A$ y de $B$.
    \end{itemize}

    \item \textbf{Pruebas:}
    \begin{itemize}
        \item a) $A \subseteq A \cup B$: Todo elemento $x \in A$ hace verdadera la proposición $x \in A \lor x \in B$.
        \item b) $A \cap B \subseteq A$: Si $x \in A \land x \in B$, entonces necesariamente $x \in A$.
    \end{itemize}

    \item \textbf{Sean $A=\{1\}, B=\{1,2\}$:}
    \begin{itemize}
        \item a) $A \in B$: \textbf{Falso}. $A$ es un subconjunto, no un elemento.
        \item b) $A \subseteq B$: \textbf{Verdadero}. El elemento 1 está en $B$.
        \item c) $1 \in A$: \textbf{Verdadero}.
    \end{itemize}

    \item \textbf{Ejercicio 3 con $A=\{1\}$ y $B=\{\{1\}, 1\}$:}
    Calcular $A - B$. El elemento 1 está en $A$ y también en $B$. Por tanto, $A - B = \emptyset$.
    
    \item \textbf{Subconjuntos de $S=\{1,2,3,4\}$:} \\
    Son $2^4 = 16$. $\emptyset, \{1\}, \{2\}, \{3\}, \{4\}, \{1,2\}, \{1,3\}, \dots, \{1,2,3,4\}$.

    \item \textbf{Problema de los 24 libros:}
    Datos: $|A|=8, |B|=13, |C|=13$. Intesecciones: $|A \cap B|=5, |A \cap C|=3, |B \cap C|=6, |A \cap B \cap C|=2$.
    Fórmula de unión:
    $|A \cup B \cup C| = 8 + 13 + 13 - (5+3+6) + 2 = 34 - 14 + 2 = 22$.
    
    \begin{enumerate}
        \item \textbf{Ninguno:} $24 - 22 = 2$ libros.
        \item \textbf{Únicamente uno:} 
        Solo A: $8 - (3+1+2) = 2$.
        Solo B: $13 - (3+4+2) = 4$.
        Solo C: $13 - (1+4+2) = 6$.
        Total = $2+4+6 = 12$ libros.
        \item \textbf{No integración (C) ni Bolzano (B):}
        Equivale a los que están solo en A o en ninguno. $|A \text{ solo}| + |\text{Ninguno}| = 2 + 2 = 4$ libros.
    \end{enumerate}
\end{enumerate}

\subsection{1.4.4 Ejercicios Propuestos con Python}

\textbf{1. Generador de Tablas de Verdad}
\begin{lstlisting}[language=Python]
print("p\tq\tr\tResult")
print("-" * 30)
for p in [True, False]:
    for q in [True, False]:
        for r in [True, False]:
            resultado = (p and q) or (not r)
            print(f"{p}\t{q}\t{r}\t{resultado}")
\end{lstlisting}

\textbf{2. Validador de Subconjuntos}
\begin{lstlisting}[language=Python]
entrada1 = input("Lista 1 (sep por comas): ").split(',')
entrada2 = input("Lista 2 (sep por comas): ").split(',')
A = set(map(int, entrada1))
B = set(map(int, entrada2))

if A.issubset(B):
    print("A es subconjunto de B")
else:
    print("A NO es subconjunto de B")
\end{lstlisting}

\textbf{4. El Problema de los 24 Libros (Script)}
\begin{lstlisting}[language=Python]
U = 24
A, B, C = 8, 13, 13
AB, AC, BC = 5, 3, 6
ABC = 2

Union = A + B + C - (AB + AC + BC) + ABC
Ninguno = U - Union
print(f"Libros que no estudian ningún tema: {Ninguno}")
\end{lstlisting}

\newpage
\section{Capítulo 2: Sistemas Numéricos}

\subsection{2.2 y 2.3 Ejercicios de Aritmética (MCD, MCM)}

\textbf{1. Cálculo de MCD y MCM}
\begin{itemize}
    \item a) 450 y 360.
    $450 = 2 \cdot 3^2 \cdot 5^2$. $360 = 2^3 \cdot 3^2 \cdot 5$.
    \textbf{MCD:} $2 \cdot 3^2 \cdot 5 = 90$.
    \textbf{MCM:} $2^3 \cdot 3^2 \cdot 5^2 = 8 \cdot 9 \cdot 25 = 1800$.
\end{itemize}

\textbf{2. Algoritmo de Euclides (Identidad de Bézout)}
\begin{itemize}
    \item a) 123 y 277.
    $277 = 2(123) + 31$ \\
    $123 = 3(31) + 30$ \\
    $31 = 1(30) + 1$ \\
    $30 = 30(1) + 0$ \\
    MCD = 1. Expresión: $1 = 31 - 30 = 31 - (123 - 3\cdot 31) = 4\cdot 31 - 123 = 4(277 - 2\cdot 123) - 123 = 4(277) - 9(123)$.
    $1 = 4(277) - 9(123)$.
\end{itemize}

\textbf{4. Problema de las Campanas}
MCM de 45, 60 y 90.
$45=3^2 \cdot 5$, $60=2^2 \cdot 3 \cdot 5$, $90=2 \cdot 3^2 \cdot 5$.
MCM = $2^2 \cdot 3^2 \cdot 5 = 4 \cdot 9 \cdot 5 = 180$ minutos.
Sonarán 3 horas después de las 8:00 a.m., es decir, a las \textbf{11:00 a.m.}

\textbf{5. Problema de Baldosas}
Se busca el MCD de 420 y 480 para el tamaño máximo.
$420 = 10 \cdot 42 = 2^2 \cdot 3 \cdot 5 \cdot 7$.
$480 = 10 \cdot 48 = 2^5 \cdot 3 \cdot 5$.
MCD = $2^2 \cdot 3 \cdot 5 = 60$.
a) Lado de la baldosa: \textbf{60 cm}.
b) Cantidad: $(420/60) \times (480/60) = 7 \times 8 = \textbf{56 baldosas}$.

\textbf{6. Ecuaciones Diofánticas (Granjero)}
$100x + 17y = 1800$.
Despejando: $17y = 100(18-x)$. Como 17 es primo y no divide a 100, debe dividir a $(18-x)$.
Posibles valores para $x$ (vacas, $x \ge 0$): $x=1$ o $x=18$.
Si $x=1$, $17y = 1700 \Rightarrow y=100$.
Si $x=18$, $17y = 0 \Rightarrow y=0$.
\textbf{Solución lógica:} 1 vaca y 100 ovejas. (O 18 vacas y ninguna oveja).

\textbf{12. Sucesión de Fibonacci (Reto)}
Propiedad: Dos números de Fibonacci consecutivos son coprimos.
$MCD(F_n, F_{n+1}) = 1$. Se demuestra aplicando Euclides: $F_{n+1} = 1 \cdot F_n + F_{n-1}$, el MCD se traslada a $(F_n, F_{n-1})$ hasta llegar a $(1,1)$.

\subsection{2.8 Ejercicios Propuestos Números Reales}

\textbf{1. Clasificación:}
\begin{itemize}
    \item $1.4142$: Racional (decimal finito).
    \item $\sqrt{2}$: Irracional.
    \item $3.1416$: Racional (decimal finito).
    \item $e$: Irracional.
    \item $0.333\dots$: Racional ($1/3$).
\end{itemize}

\textbf{4. Inverso Multiplicativo:}
\begin{itemize}
    \item a) De 5: $1/5$.
    \item b) De $3/4$: $4/3$.
    \item c) De $0.5$ ($1/2$): $2$.
\end{itemize}

\subsection{2.9 Laboratorio con Python}

\textbf{Ejercicio Propuesto 2 (Ecuación)}
\begin{lstlisting}[language=Python]
x = 3
resultado = (x**4 == -1)
print(f"¿Es solución? {resultado}") 
# Salida esperada: False, pues 81 != -1
\end{lstlisting}

\textbf{Ejercicio Propuesto 3 (Múltiplos)}
\begin{lstlisting}[language=Python]
# Múltiplos de 3 del 3 al 15
multiplos = list(range(3, 16, 3))
print(multiplos) # [3, 6, 9, 12, 15]
\end{lstlisting}

\subsection{2.10 Ejercicios de Profundización (MCD Python)}

\textbf{Código para MCD Automatizado:}
\begin{lstlisting}[language=Python]
def mcd_euclides(a, b):
    while b != 0:
        a, b = b, a % b
    return a

num1 = 123456
num2 = 7890
print(f"El MCD es: {mcd_euclides(num1, num2)}")
\end{lstlisting}

\end{document}
